\documentclass[11pt]{article}
\usepackage[a4paper,margin=1in]{geometry}
\usepackage{amsmath,amssymb,amsthm}
\usepackage{booktabs}
\usepackage{hyperref}
\usepackage{enumitem}

\newtheorem{theorem}{Theorem}
\newtheorem{lemma}[theorem]{Lemma}
\newtheorem{proposition}[theorem]{Proposition}
\newtheorem{corollary}[theorem]{Corollary}
\newtheorem{remark}[theorem]{Remark}

\title{The Vertex-Minimal Nontrivial Nucleus: An Exhaustive Computational Classification}
\author{Bergfinnur Durhuus and Marta Pavelka}
\date{}

\begin{document}
\maketitle

\noindent
Supplementary computational details for
\emph{Counting Triangulated 3-Spheres, Nuclei, and Polyhedra}.


\section{Nuclei and the search bounds}

A \emph{nucleus} is a finite simplicial triangulation $N$ of the $3$-ball such that
\begin{enumerate}[label=(\roman*)]
    \item every vertex of $N$ lies on $\partial N$, and
    \item every interior triangle contains at most one boundary edge.
\end{enumerate}
It is \emph{nontrivial} if it has more than one tetrahedron.  Write
$v=f_0(N)$ and $t=f_3(N)$.  An edge is \emph{interior} if it does not
belong to $\partial N$; since every vertex belongs to the boundary, interior
edges are also called spanning edges.

The classification below uses one geometric input, proved independently in
the structural part of the companion paper:
\begin{quote}
\emph{All-interior-edge tetrahedron theorem.} Every nontrivial nucleus
contains a tetrahedron all six of whose edges are interior.
\end{quote}
This result is not inferred from the enumeration.  It provides a necessary
constraint for the search; none of the arguments below requires the converse.

\begin{lemma}\label{lem:interior-edges}
For a triangulated $3$-ball $N$ whose vertices all lie on its boundary,
\[
 f_1^{\mathrm{int}}(N)=t-v+3.
\]
\end{lemma}

\begin{proof}
The boundary is a triangulated $2$-sphere on $v$ vertices, so
\[
 f_1(\partial N)=3v-6,\qquad f_2(\partial N)=2v-4.
\]
Counting tetrahedron--triangle incidences gives
\[
 4t=2f_2^{\mathrm{int}}(N)+2v-4,
 \qquad f_2(N)=2t+v-2.
\]
Euler's formula $v-f_1(N)+f_2(N)-t=1$ now yields
$f_1(N)=t+2v-3$.  Subtracting the $3v-6$ boundary edges proves the claim.
\end{proof}

\begin{proposition}\label{prop:tbound}
If $N$ is a triangulated $3$-ball with all vertices on its boundary, then
\[
 \boxed{\ t\le\binom v2-2v+3.\ }
\]
In particular, $t\le28$ for $v\le10$, and $t\le36$ for $v=11$.
\end{proposition}

\begin{proof}
At most $\binom v2$ edges occur in $N$, of which $3v-6$ lie on
the boundary.  Lemma~\ref{lem:interior-edges} thus gives
\[
 t-v+3=f_1^{\mathrm{int}}(N)
 \le\binom v2-(3v-6),
\]
which rearranges to the stated bound.  Its right-hand side is increasing
for $v\ge3$ and takes the values $28$ and $36$ at $v=10$ and $v=11$.
\end{proof}

For $v=11$, the all-interior-edge tetrahedron theorem also gives
$f_1^{\mathrm{int}}(N)\ge6$.  Lemma~\ref{lem:interior-edges} then implies
\begin{equation}\label{eq:t-range}
 14\le t\le36.
\end{equation}

\section{Enumeration of the possible boundaries}

The boundary of a nucleus is a simplicial triangulation of the $2$-sphere
on the same vertex set.  Its combinatorial types are generated using
\texttt{plantri}, developed by Gunnar Brinkmann and Brendan McKay.
A maximal planar graph with its planar embedding determines a simplicial
triangulated $2$-sphere.  For at least four vertices the embedding is unique
up to reflection, so the triangulation mode of \texttt{plantri} lists the
possible boundary spheres once up to isomorphism.

The numbers of boundary types are
\[
\begin{array}{c|rrrrrrrr}
 v&4&5&6&7&8&9&10&11\\ \hline
 \#S&1&1&2&5&14&50&233&1249.
\end{array}
\]
Thus there are $306$ types for $4\le v\le10$ and $1249$ types for $v=11$.
The lists used in this classification are included in the reproducibility
package as \texttt{data/boundaries\_v4.asc} through
\texttt{data/boundaries\_v11.asc}. The \texttt{plantri} program can be
obtained separately from its official distribution at
\url{https://users.cecs.anu.edu.au/~bdm/plantri/}.
Different software versions may list the same combinatorial types in a
different order; the original input files are therefore retained.

\section{A necessary binary filling model}\label{sec:binary-model}

Fix a boundary sphere $S$ on vertex set $V$, and write $E(S)$ for its
edges.  Any triangle $F\notin S$ containing two or three edges of $E(S)$
is \emph{forbidden}: if it belonged to a selected tetrahedron, the prescribed
boundary incidences would force it to be an interior triangle, contrary to
the nucleus condition.  We therefore discard every candidate tetrahedron
containing a forbidden triangle.

For each remaining four-element set $Q\subseteq V$ introduce a binary
variable $x_Q$, recording whether $Q$ is selected as a tetrahedron.  For
each nonboundary, nonforbidden triangle $F$ introduce a binary variable
$y_F$, recording whether it is an interior triangle.  For each nonboundary
edge $e$ introduce a binary variable $z_e$, recording whether it occurs in
the filling.  The model imposes the following necessary conditions.

\begin{enumerate}[label=(\arabic*),leftmargin=*]
\item \textbf{Boundary triangles.} Every triangle of $S$ is incident with
exactly one selected tetrahedron:
\[
 \sum_{Q\supset F}x_Q=1\qquad(F\in S).
\]

\item \textbf{Nonboundary triangles.} Every allowed triangle outside $S$
is incident with zero or two selected tetrahedra:
\[
 \sum_{Q\supset F}x_Q=2y_F
 \qquad(F\notin S\text{ allowed}).
\]
Forbidden triangles occur in no candidate tetrahedron.

\item \textbf{Nonboundary edges.} For every $e\notin E(S)$, with
$Q$ ranging over candidate tetrahedra containing $e$, impose
\[
 x_Q\le z_e\quad(e\subset Q),\qquad
 z_e\le\sum_{Q\supset e}x_Q.
\]
Consequently $z_e=1$ if and only if $e$ occurs in a selected tetrahedron.

\item \textbf{Tetrahedron bound.} Proposition~\ref{prop:tbound} yields
\[
 \sum_Qx_Q\le\binom v2-2v+3.
\]

\item \textbf{An all-interior-edge tetrahedron.} The independent structural
theorem gives
\[
 \sum_{Q:\,\binom Q2\cap E(S)=\varnothing}x_Q\ge1.
\]
If no candidate $Q$ has all six edges outside $E(S)$, the system is
immediately infeasible.

\item \textbf{Vertex-link Euler characteristics.} Every vertex link in
a triangulated $3$-ball with all vertices on the boundary is a disk and
has Euler characteristic one.  The vertices, edges, and triangles of
$\operatorname{lk}_N(a)$ correspond to the edges, triangles, and tetrahedra
of $N$ containing $a$.  Thus, for each $a\in V$,
\[
 \#\{e\in E(S):a\in e\}
 +\sum_{\substack{e\notin E(S)\\a\in e}}z_e
 -\#\{F\in S:a\in F\}
 -\sum_{\substack{F\notin S\text{ allowed}\\a\in F}}y_F
 +\sum_{Q\ni a}x_Q=1.
\]
\end{enumerate}

These conditions are \emph{necessary}, not asserted to be sufficient.
Every nontrivial nucleus with boundary $S$ determines a feasible binary
assignment by taking the indicators of its tetrahedra, interior triangles,
and interior edges.  Conversely, a feasible assignment may fail to triangulate
a $3$-ball.  Notice also that once the tetrahedron indicators $x_Q$ are
fixed, the triangle and edge equations uniquely determine all $y_F$ and
$z_e$.  This will justify blocking solutions using only the $x_Q$ variables.

\section{The lower bound of eleven vertices}

For each of the $306$ boundary spheres with $4\le v\le10$, the model in
Section~\ref{sec:binary-model} was solved as a binary feasibility problem
using SciPy's interface to HiGHS.  All $306$ instances were reported
infeasible.  The boundary files and machine-readable feasibility results
are included in the reproducibility package.

\begin{theorem}\label{thm:11vertices}
Every nontrivial nucleus has at least $11$ vertices.
\end{theorem}

\begin{proof}
Suppose a nontrivial nucleus $N$ has $v\le10$ vertices, and put
$S=\partial N$.  Up to relabeling, $S$ is one of the $306$ boundary spheres
from Section~2.  The tetrahedra, interior triangles, and interior edges of
$N$ determine a feasible assignment for the corresponding binary model,
by the observations in Section~\ref{sec:binary-model}.  But the exhaustive
computation reports all such models infeasible, a contradiction.
\end{proof}

The argument does not need $3$-ball recognition: infeasibility already holds
for a system of necessary combinatorial constraints.

\section{Classification on eleven vertices}

There are $1249$ possible boundary types, and by~\eqref{eq:t-range} it
suffices to consider $14\le t\le36$.  For each boundary $S$ and each $t$
in that range, the model is augmented by
\begin{equation}\label{eq:exact-t}
 \sum_Qx_Q=t.
\end{equation}

Whenever the solver finds a feasible tetrahedron indicator vector
$\bar x\in\{0,1\}^m$, the program appends the \emph{no-good} inequality
\[
 \sum_{i:\bar x_i=1}x_i-\sum_{i:\bar x_i=0}x_i
 \le |\{i:\bar x_i=1\}|-1.
\]
This excludes exactly that tetrahedron set.  As noted above, the $y$ and
$z$ variables are uniquely determined by $x$, so no different filling is
lost.  Iteration continues until HiGHS reports infeasibility; hence the
procedure enumerates all feasible binary fillings, subject to the solver's
reported decisions.

For each enumerated filling, the program checks the exact simplicial
incidences and local links and computes homology over $\mathbb F_2$.
Every $3$-ball has Betti numbers $(1,0,0,0)$, so any other Betti vector
excludes the filling.  The reported exhaustive counts are:

\begin{center}
\begin{tabular}{@{}c r r@{}}
\toprule
$t$ & feasible binary fillings & nuclei\\
\midrule
$14$--$29$ & $0$ & $0$\\
$30$ & $8$ & $0$\\
$31$ & $54$ & $0$\\
$32$ & $213$ & $0$\\
$33$ & $504$ & $0$\\
$34$ & $623$ & $0$\\
$35$ & $360$ & $1$\\
$36$ & $67$ & $0$\\
\midrule
Total & $1829$ & $1$\\
\bottomrule
\end{tabular}
\end{center}

Of the $1829$ feasible fillings, $1828$ have mod-$2$ Betti numbers
$(1,1,1,0)$ and therefore are not $3$-balls.  The only remaining filling
occurs at $t=35$, on the stored boundary representative numbered $689$
(using the program's zero-based indexing).  Its $35$ facets are recorded in
$\texttt{results/n00\_facets.json}$.  We call it $N_{\min}$; it is the
complex denoted \texttt{N00} in the computation.  Its vertices are labeled
$0,\ldots,10$ in the data, corresponding to $1,\ldots,11$ in the main paper.
A direct computation gives
\[
 f(N_{\min})=(11,54,79,35),\qquad
 f(\partial N_{\min})=(11,27,18).
\]

The positive verification is separate from the homology rejection of the
other fillings.  The facet list passes exact checks of the prescribed
spherical boundary, triangle incidences, reducedness, dual connectivity,
vertex-disk links, and the appropriate interval or circle edge links.
In addition, a certificate of $89$ elementary simplicial collapses to a
vertex is supplied and independently replayed.  The stricter replay
verifier also checks that a supposedly free face has no additional
cofaces in higher dimensions.  These checks show that the complex is a
contractible PL $3$-manifold with boundary $S^2$ and hence a PL $3$-ball.
Thus it is indeed a nucleus.

\begin{theorem}[Classification at the minimum]\label{thm:unique11}
Up to simplicial isomorphism, $N_{\min}$ is the unique nontrivial nucleus
with $11$ vertices.  It has exactly $35$ tetrahedra and
\[
 f(N_{\min})=(11,54,79,35).
\]
In particular, there is no nontrivial $11$-vertex nucleus with $36$
tetrahedra.
\end{theorem}

\begin{proof}
Let $N$ be any nontrivial $11$-vertex nucleus.  Its boundary agrees,
after relabeling, with one of the $1249$ generated spheres, and its
number of tetrahedra satisfies $14\le t\le36$.  Therefore its tetrahedra
induce one of the feasible binary fillings exhaustively enumerated for
that boundary and $t$.  Every filling other than the displayed $N_{\min}$
has nonzero first and second homology over $\mathbb F_2$, whereas
$N_{\min}$ is positively verified to be a nucleus.  This proves uniqueness
and the stated facet count.
\end{proof}

\begin{corollary}\label{cor:smallest}
The minimum number of vertices of a nontrivial nucleus is $11$, and
there is exactly one combinatorial type at this minimum: $N_{\min}$.
\end{corollary}

\begin{proof}
Theorem~\ref{thm:11vertices} rules out fewer vertices, while
Theorem~\ref{thm:unique11} establishes existence and uniqueness at $11$.
\end{proof}

\section{Reproducibility and scope of the certificates}

The reproducibility package contains the stored \texttt{plantri} outputs,
not the program's source, which is available separately from its official
website. It contains the common binary-model implementation, the $v=11$ classification driver,
the $1249$ stored boundary representatives, the recorded $v=11$ count table,
the facet list and collapse certificate for $N_{\min}$, and their checksums.
It also contains the boundary data and complete
machine-readable results for the $v\le10$ lower-bound run, together with a
verifier of the collapse certificate.

From the package root, the principal commands are:
\begin{verbatim}
# Obtain and compile plantri from its official website first.
PLANTRI=/absolute/path/to/plantri
python src/nucleus_lower_bound.py \
    --plantri "$PLANTRI" \
    --min-v 4 --max-v 10 --out reproduced_lower_bound
python src/classify_v11.py \
    --plantri "$PLANTRI" \
    --out reproduced_v11.json
python src/verify_n00_collapse_strict.py \
    results/n00_collapse_certificate.json
sha256sum -c SHA256SUMS.txt
\end{verbatim}

\begin{remark}[Computational status]\label{rem:computational-status}
The $v\le10$ infeasibility results and the positive $N_{\min}$
checks are reproducible by the supplied programs. The $v=11$ counts
are provided as machine-readable summaries of the exhaustive enumeration;
they can be recomputed with the included driver. The negative results
do not include independently checkable solver proof objects.

The negative conclusions depend on the correctness of \texttt{plantri}, the
binary-model encoding, and the HiGHS feasibility and infeasibility decisions.
The package does not contain independently checkable certificates of
infeasibility for each MILP instance.  In contrast, the positive collapse
certificate is replayable independently of the MILP solver.  These
qualifications apply to the computational classification theorem above.
\end{remark}

\end{document}
